% =============================================================
%  ch24_portalsystem.tex  --  Mechanischer Aufbau: Portalsystem (Gantry)
% =============================================================

\chapter{Mechanischer Aufbau: Portalsystem (Gantry)}

\section{Grundkonzept}
Statt eines Delta-Roboters ein \textbf{kartesisches Portalsystem}: große
Linearachsen positionieren den Schläger grob im Raum, ein kleiner Arm am Ende
führt die Schlagbewegung aus und stellt den Schlägerwinkel. Bezug
Tischtennisplatte (ITTF): $2{,}74\,\text{m} \times 1{,}525\,\text{m}$, Höhe
$0{,}76\,\text{m}$, Netz $15{,}25\,\text{cm}$. Die \textbf{X-Achse} überspannt
die Plattenbreite ($1{,}525\,\text{m}$).

\section{Geschachtelte Achsen}
\begin{itemize}
\item \textbf{X-Achse} (Plattenbreite, horizontal): zwei parallele
  Führungsstäbe, Zahnriemenantrieb, X-Schlitten. Bauprinzip wie die
  Bambu-Lab-P2S-X-Achse -- \emph{glatte Hohlstahl-Stäbe + GT2-Zahnriemen},
  leichter Schlitten.
\item \textbf{Z-Achse} (Höhe, vertikal): nach demselben Prinzip, montiert
  \emph{im} X-Schlitten. Der Abstand der X-Führungsstäbe definiert den Bauraum
  und damit den \textbf{Hub der Z-Achse}.
\item \textbf{Y-Achse} (Tiefe): entweder starr (fester Ausleger) oder als
  dritte Linearachse auf dem Z-Schlitten.
\item \textbf{Schlagarm}: auf dem Y/Z-Schlitten ein kleiner Roboterarm mit
  Schlägereinspannung.
\end{itemize}

\begin{warnbox}[title={Zahnriemen, nicht Keilriemen}]
Für \emph{Positionierung} brauchst du einen \textbf{Zahnriemen}
(GT2-Zahnriemen), keinen Keilriemen (V-Riemen). Ein Keilriemen überträgt per
Reibschluss und \emph{schlupft} -- damit verlierst du Positionsgenauigkeit.
Zahnriemen sind formschlüssig, schlupffrei und der Standard in 3D-Druckern/Plottern.
\end{warnbox}

\section{Antriebstopologie und Hardware-Anschluss}
\begin{longtable}{@{}p{3.2cm}p{11.4cm}@{}}
\toprule
\textbf{Element} & \textbf{Umsetzung} \\
\midrule
X-/Z-/Y-Motoren  & NEMA-17, möglichst \emph{rahmenfest} (Riemen treibt Schlitten) \\
Treiber          & TMC2209 (UART) je Achse, StallGuard als Last-/Endlagenhinweis \\
Endschalter      & optische/mechanische Endstops je Achse zum Homing \\
Strom/Spannung   & INA226-Modul je Treiberzweig (I\textsuperscript{2}C), optional \\
Controller       & ESP32 (micro-ROS): bündelt STEP/DIR bzw.\ UART, liest Endstops + INA226 \\
Versorgung       & 24\,V Netzteil für Motoren, getrennte 5\,V-Logik; Sternmasse \\
Schlagarm-Motoren & klein, idealerweise basisnah (Riemen-/Seilzug zu distalen Gelenken) \\
\bottomrule
\end{longtable}

\begin{center}
\begin{tikzpicture}[scale=1.0,every node/.style={font=\scriptsize}]
% frame
\draw[thick,accent] (0,0) rectangle (7,4.2);
\node[accent] at (3.5,4.5) {Rahmen (rahmenfest, X-Motor + Idler)};
% X rails
\draw[gray,line width=2pt] (0.5,3.6)--(6.5,3.6);
\draw[gray,line width=2pt] (0.5,2.4)--(6.5,2.4);
\node[gray] at (-0.2,3.0) {X};
% X carriage holding Z
\draw[fill=accent!12,draw=accent] (2.6,2.2)rectangle(4.4,3.8);
\node[accent] at (3.5,3.95) {X-Schlitten};
% Z rails inside carriage
\draw[accent2,line width=1.5pt] (2.9,2.3)--(2.9,3.7);
\draw[accent2,line width=1.5pt] (4.1,2.3)--(4.1,3.7);
\node[accent2] at (3.5,3.0) {Z};
% Z carriage
\draw[fill=accent2!12,draw=accent2] (2.95,2.55)rectangle(4.05,3.0);
\node[accent2] at (3.5,2.78) {Z-Schlitten};
% arm
\draw[line width=1.4pt,black] (3.5,2.55)--(3.5,1.0);
\draw[fill=black!8] (3.2,0.6) rectangle (3.8,1.0);
\node at (3.5,0.8) {Arm};
\draw[fill=accent!20,draw=accent] (3.35,0.2) rectangle (3.65,0.6);
\node[rotate=90] at (3.5,0.4) {\tiny Schl\"ager};
% lever arm annotation
\draw[<->,accent2] (4.0,2.55)--(4.0,0.4);
\node[accent2,right] at (4.05,1.5) {Hebelarm $L$};
\end{tikzpicture}
\captionof{figure}{Schematischer Frontblick: X-Schlitten trägt die im
Rod-Abstand geschachtelte Z-Achse; der Schlagarm hängt als Ausleger (Hebelarm
$L$) darunter.}
\end{center}
